%% starter: A lab report in APS style, with references.bib and a figures folder
\documentclass[reprint,aps]{revtex4-2}
\usepackage{amsmath}
\usepackage{graphicx}
\usepackage{siunitx}

\begin{document}

\title{{{Title}}}
\author{{{Author}}}
\affiliation{Department of Physics}
\date{\today}

\begin{abstract}
What you measured, how, and what you found, in a few sentences.
\end{abstract}

\maketitle

\section{Introduction}
Why the experiment is worth doing, and what it sets out to measure~\cite{example}.

\section{Theory}
The period of a simple pendulum of length $L$ is
\begin{equation}
  T = 2\pi\sqrt{\frac{L}{g}},
  \label{eq:period}
\end{equation}
so a plot of $T$ against $\sqrt{L}$ is a straight line of slope $2\pi/\sqrt{g}$.

\section{Methods}
The apparatus and procedure. The pendulum was \qty{0.50}{m} long.

\section{Results}
Fig.~\ref{fig:data} shows the data, and fitting Eq.~\eqref{eq:period} to it gives $g$.

\begin{figure}
  % Put your plot in figures/ and uncomment the next line, naming it.
  % \includegraphics[width=\linewidth]{figures/plot.pdf}
  \fbox{\parbox[c][4cm][c]{0.9\linewidth}{\centering Your figure here}}
  \caption{The period $T$ against $\sqrt{L}$, with a straight-line fit.}
  \label{fig:data}
\end{figure}

\section{Conclusion}
What the results show, and how they compare with the accepted value.

\begin{acknowledgments}
Thanks to your lab partner and demonstrator.
\end{acknowledgments}

\bibliography{references}

\end{document}
